% !TeX program = xelatex
% !TeX program = xelatex
% ===========================================================================
%  tech_common.tex -- 两份技术文档共用的导言区
% ===========================================================================
\documentclass[11pt,a4paper]{ctexart}

\usepackage[top=2.4cm,bottom=2.4cm,left=2.4cm,right=2.4cm]{geometry}
\usepackage{amsmath,amssymb,amsthm}
\usepackage{graphicx}
\usepackage{booktabs}
\usepackage{longtable}
\usepackage{array}
\usepackage{enumitem}
\usepackage{xcolor}
\usepackage{listings}
\usepackage{tcolorbox}
\usepackage{caption}
\usepackage{fancyhdr}
\usepackage{titlesec}
\usepackage{hyperref}

\hypersetup{colorlinks=true,linkcolor=blue!55!black,citecolor=blue!55!black,
            urlcolor=blue!55!black,bookmarksnumbered=true}

% ---------------------------------------------------------------- 字体/版式
\setlength{\emergencystretch}{2.5em}
\setlength{\parskip}{0.35em}
\setlength{\parindent}{2em}
\linespread{1.12}

% ---------------------------------------------------------------- 代码样式
\definecolor{codebg}{RGB}{247,247,250}
\definecolor{codekw}{RGB}{0,90,160}
\definecolor{codecm}{RGB}{0,128,96}
\definecolor{codest}{RGB}{170,60,60}
\lstset{
  basicstyle=\ttfamily\footnotesize,
  backgroundcolor=\color{codebg},
  frame=single, framerule=0.4pt, rulecolor=\color{gray!45},
  breaklines=true, breakatwhitespace=false,
  numbers=left, numberstyle=\tiny\color{gray!80}, numbersep=6pt,
  keywordstyle=\color{codekw}\bfseries, commentstyle=\color{codecm}\itshape,
  stringstyle=\color{codest}, showstringspaces=false, columns=flexible,
  tabsize=2, extendedchars=true, captionpos=b,
  xleftmargin=1.6em, framexleftmargin=1.4em
}
\renewcommand{\lstlistingname}{代码}

% ---------------------------------------------------------------- 提示框
\newtcolorbox{keybox}[1][]{colback=blue!4,colframe=blue!45!black,
  boxrule=0.6pt,arc=2pt,left=5pt,right=5pt,top=4pt,bottom=4pt,
  fonttitle=\bfseries,#1}
\newtcolorbox{warnbox}[1][]{colback=orange!6,colframe=orange!55!black,
  boxrule=0.6pt,arc=2pt,left=5pt,right=5pt,top=4pt,bottom=4pt,
  fonttitle=\bfseries,#1}

% ---------------------------------------------------------------- 定理环境
\newtheorem{theorem}{定理}
\newtheorem{proposition}{命题}
\newtheorem{lemma}{引理}
\renewcommand{\proofname}{\heiti 证明}

% ---------------------------------------------------------------- 页眉页脚
\pagestyle{fancy}
\fancyhf{}
\fancyhead[L]{\small\color{gray!70}\leftmark}
\fancyhead[R]{\small\color{gray!70}\thepage}
\renewcommand{\headrulewidth}{0.3pt}
\renewcommand{\headrule}{\hbox to\headwidth{\color{gray!40}\leaders\hrule height \headrulewidth\hfill}}

\titleformat{\section}{\normalfont\heiti\Large}{\thesection}{0.6em}{}
\titleformat{\subsection}{\normalfont\heiti\large}{\thesubsection}{0.6em}{}
\titleformat{\subsubsection}{\normalfont\heiti\normalsize}{\thesubsubsection}{0.6em}{}

\captionsetup{font=small,labelfont=bf,labelsep=quad}


\title{\heiti 无线电干扰源定位与清除\\[6pt]
       \large 建模思路、算法设计与实测依据\\[4pt]
       \normalsize 2026 年全国大学生数学建模竞赛 B 题}
\author{}
\date{\today}

\begin{document}
\maketitle
\thispagestyle{fancy}

\begin{abstract}
\noindent
本文围绕无线电干扰源的快速自动定位与清除，建立了几何定位、站位优化、
在线搜索规划与可证终止四层模型，并按附件协议实现本地模拟环境完成验证。

问题一把带正负 1 度误差的示向度解释为以检测点为顶点的角楔，定位区域取各楔之交，
证明其为凸多边形，给出 $O(n^2)$ 顶点枚举算法与基于回收锥的有界性判据，
用旋转卡壳求直径、Welzl 算法求最小包围圆，并给出一阶解析式。
问题二在含正负 1 度误差的最小包围圆意义下求解第二观测点，得到 b*=1004.1 m，
按该规则观测时直径上限可达 134.0 m。
问题三用七点覆盖定理与滚动时域路径规划，60 组演练完成率 1.0000，平均 326.8 秒。
问题四给出定向源的认证判据、射线逼近引理与三角格点间距条件，
在四个独立算例池共 120 例上最差完成率 0.9978，平均 595.9 秒每源。

本文的重点之一是把不可达的目标也如实界定。经逐步核算，一局必须访问约 58 个位置点，
这批点的最优巡回约 20000 m，折合 320 秒每源，已接近 300 至 400 秒的目标区间下界；
加上检测开销后现实下界在 450 秒每源以上。研究过程中设计并实测了 9 项结构性替代方案
与 30 余项参数尝试，其中大部分为负结果，全部记录在案。
\end{abstract}

\tableofcontents
\newpage

% =====================================================================
\section{问题与约束}

目标区域内存在 10 至 16 个无线电干扰源，占用 20 个信道中的若干互不相同的信道，
工作频段均在测向机接收范围内。干扰源分为两类：全向源向四周辐射，
定向源向某个 180 度扇区辐射。每个源的有效接收半径为 1000 至 1500 m 之间的某一定值，
互不相同。信道号已知，但每个信道是否有源、源的位置、辐射类型与朝向均未知。

搭载测向机的机器人在半径 1800 m 的圆形区域内移动，速度 5 m/s。
测向机同一时刻只能接收一个信道，切换信道耗时 1 秒，测量一次耗时 5 秒。
测量给出的是从当前位置指向该源的真实方位角，误差为不超过正负 1 度。
按附件说明，同一位置对同一信道的测量误差固定，即误差是位置的函数，
重复测量同一点不会带来新信息。

机器人在距离干扰源 20 m 以内时可以将其清除，清除动作成功耗时 5 秒，
失败 3 秒。单局虚拟时间上限 360000 秒，真实时间上限 1200 秒。

这两条约束决定了整套算法的形状：测量成本高，误差不可通过重复测量消除，
清除需要物理接近，因此问题的核心是信息效率与行程效率的联合优化。

% =====================================================================
\section{建模总框架}

四个子问题共享同一套几何基础，差别在于信息结构与代价结构：

\begin{center}
\begin{tabular}{llll}
\toprule
问题 & 信息结构 & 代价结构 & 求解目标 \\
\midrule
一 & 已知各次测量位置 & 无 & 确定定位区域与不确定度 \\
二 & 可自选两个测量点 & 到达距离 & 最小化最坏情形直径 \\
三 & 全向源，在线探测 & 移动与测量 & 最短时间完成清除 \\
四 & 含定向源，在线探测 & 移动与测量 & 同上，且需证明无遗漏 \\
\bottomrule
\end{tabular}
\end{center}

贯穿四问的一条主线是终止条件。在线搜索必须回答何时可以停止，
而停止的依据只能是某种可验证的判据。问题三与问题四的终止判据分别建立在
覆盖定理与凸包判据之上，这也是两问在算法上分野的根源。

% =====================================================================
\section{问题一：几何定位}

\subsection{角楔模型}

一次测量给出真实方位角落在测量值正负 1 度之内。以检测点为顶点、
以测量方位为中心、张角 2 度的角楔区域，就是该次测量所确定的位置约束。
多次测量的定位区域即所有角楔之交。

\subsection{凸性与顶点枚举}

角楔是凸集的交，因此定位区域是凸多边形。取每条边的两条边界射线，
两两求交得到候选顶点，保留同时落在所有角楔内的点，即得多边形顶点集。
算法复杂度为 $O(n^2)$ 次求交与 $O(n^3)$ 次判定。

\subsection{有界性}

区域可能有界也可能无界，取决于各角楔方向是否覆盖整个圆周。
判据是这些方向的回收锥是否为整个平面：若存在一条射线方向，
使得沿该方向可以无限远离且始终落在所有角楔内，则区域无界。
实践中检查各角楔方向的最左与最右边界射线是否留下空隙即可。

\subsection{不确定度度量}

用最小包围圆的半径作为不确定度。求法为 Welzl 随机增量算法，
期望线性时间。直径用旋转卡壳法在凸多边形上线性求得。
在两条测量读数相距较远、交会角为 gamma 的情形下，
直径有渐近式

\begin{center}
$D \approx \dfrac{\sqrt{w_1^2 + w_2^2 + 2 w_1 w_2 \cos\gamma}}{\sin\gamma}$,
\qquad $w_i = 2\varepsilon d_i$
\end{center}

其中 epsilon 为测角误差，$d_i$ 为观测点到源的距离。该式说明直径同时正比于
观测距离与误差，且反比于交会角的正弦，这为问题二的站位设计提供了直接的依据。

% =====================================================================
\section{问题二：第二观测点}

\subsection{极小极大形式}

在第一条方位已知的条件下，选择第二观测点 S2 使最坏情形直径最小。
目标函数取最小包围圆半径的三倍作为直径上界，
约束为 S2 到目标的距离不小于 $R_m$in，以保证仍然可以接收到信号。

求解结果显示最优站位近似满足 S2 到目标的距离等于 $R_m$in，
且偏离已知方位的角度约为 37 度。此时 b*=1004.1 m，最坏直径 J*=134.0 m。

\subsection{与基线的对比}

\begin{center}
\begin{tabular}{lc}
\toprule
策略 & 最坏直径 m \\
\midrule
本文极小极大站位 & 134.0 \\
均值站位 & 55.1 \\
垂直平分线站位 & 96.6 \\
最远可达站位 & 2508.7 \\
就近站位 & 1285.0 \\
在估计点原地加测 & 14.8（信息量为零） \\
\bottomrule
\end{tabular}
\end{center}

最后一行说明一个容易忽略的事实：在原地重复测量不会缩小误差，
因为误差是位置的函数。这条性质在问题三与问题四的算法设计中反复用到。

% =====================================================================
\section{问题三：全向源的搜索与清除}

\subsection{终止判据}

对全向源，若在某个位置 p 测量信道 c 得到无信号，且 p 到候选位置 q 的距离
不超过 $R_m$in，则 q 处不可能存在该源。因此收集到的无信号读数构成一组覆盖约束，
当整个目标域被这些读数覆盖时，可以证明不存在遗漏。

\subsection{七点覆盖定理}

用半径 rho 的圆周上均布六点加圆心共七个站位，其接收圆盘的覆盖半径
不超过 $R_m$in = 1000 m 时，即可覆盖整个目标域。这给出 rho 的取值上界。
本文取 rho = 1250 m 留出余量。

\subsection{在线搜索}

实际执行中，站位按滚动时域方式规划：已定位的源进入清除任务，
未解算的信道进入补测任务，两类任务混合成一张任务表，
用最近邻构造加 2-opt 改进排序，每个任务执行后重建。

\subsection{结果}

60 组演练完成率 1.0000，平均 326.8 秒每源，平均行程 13035 m。
三次正式测试分别清除 15/15、12/12、10/10。

% =====================================================================
\section{问题四：定向源的搜索与清除}

\subsection{单次无信号不足以排除}

定向源只向 180 度扇区辐射，因此在某点无信号只说明该点落在背光侧，
不能排除该处存在源。这是问题四与问题三的根本差别。

\subsection{认证判据}

\begin{keybox}[title=定理]
候选位置 q 处不可能存在定向源，当且仅当 q 落在集合
$\{\,p_j - q \;:\; |p_j - q| \le R\,\}$ 的凸包内部，
其中 $p_j$ 为所有无信号读数的位置，$R$ 为认证半径。
\end{keybox}

直观含义是：若读数从各个方向包围了 q，则无论源朝哪个方向辐射，
至少有一个读数落在受光侧，因此必然会被收到。

\subsection{格点间距条件}

对间距为 s 的三角格点，任一候选点所在三角形的三个顶点到该点的距离
不超过 s。要求这三个顶点都落在认证半径 R 内且包围该点，
即得条件 s 不超过 R。认证半径由测量半径减去格点容差得到，
数值上约为 958 m。全域所需站位数为 $\pi r_A^2 / (\frac{\sqrt3}{2} s^2)$，
代入得 12.8，取整为 13。

\subsection{射线逼近引理与贴边源}

\begin{keybox}[title=引理]
若在点 p 听到某定向源，则从 p 出发朝向该源的线段整体落在源的受光半平面内。
\end{keybox}

该引理给出了一条保底的接近路径，但只保证朝向源的那一段。
越过源之后不再有保证，这一点在实验中造成了明确的故障。

贴边且朝外辐射的源是问题四最难的情形。设源位于半径 $r_q$ 处、
沿朝外径向辐射，则观测点 p 能被收到需要满足 $(p-q)\cdot \hat u \ge 0$，
在径向上即要求 $|p| \ge r_q$。也就是说观测点必须比源更靠外。
这类源在域内的受光区是一条贴着边界的月牙，实测中最窄的只有 27 m 宽，
面积约占全域的 0.11%。

\subsection{贴边环}

基于上述几何事实，在半径 1750 m 处预先布置 12 个站位。
这些站位一方面覆盖了认证搜索中难度最高、代价最大的边界候选点，
另一方面提高了贴边源的发现率。实测把平均时间由 880.7 秒降到 595.9 秒，
完成率不变。

\subsection{结果}

四个独立算例池各 30 例，共 120 例：

\begin{center}
\begin{tabular}{lccc}
\toprule
算例池 & 最差完成率 & 平均时间 s/源 & 失败算例 \\
\midrule
A（种子 62000+） & 1.0000 & 593.4 & 0 \\
B（种子 51000+） & 0.9978 & 629.6 & 1 \\
C（种子 12000+） & 1.0000 & 580.1 & 0 \\
D（种子 30000+） & 1.0000 & 579.9 & 0 \\
\midrule
合计 & 0.9978 & 595.9 & 1/120 \\
\bottomrule
\end{tabular}
\end{center}

在另取的 60 例全新池上共 800 个源，漏源 2 个，占 0.250\%。
两个漏源都是贴边朝外辐射，与朝外径向的夹角余弦分别为 0.980 与 0.910，
域内受光区分别占 0.28\% 与 3.54\%。这类构型需要沿 11.3 km 边界
每隔 25 m 布设站位才能保证覆盖，约需 450 个站位，时间预算不允许。

% =====================================================================
\section{模拟器与验证}

由于官方模拟器无法获取，本文按附件协议实现了本地模拟环境，
并与开源项目 Jammers\-Simulator\-Tool 做了交叉核对：12 个常量全部一致，
300 次随机动作的代价模型偏差小于 $10^{-9}$ 秒。
唯一分歧是该开源项目的离线桩件每次调用都重新抽取示向度误差，
违反附件关于同一位置误差固定的说明，本文按附件实现。

验证套件共 68 项，分为八组：协议 18 项、计时 8 项、物理 16 项、
响应契约 8 项、交叉核对 5 项、时间守恒 4 项、增量等价 5 项、策略保证 4 项。
全部通过。其中增量等价组专门核对认证扫描的增量实现与暴力实现结果一致，
该组在开发过程中抓到过一个把 1253 个未覆盖点截断为 1 个的严重缺陷。

% =====================================================================
\section{性能分析：时间花在哪里}

对最终配置做 10 例逐腿审计：

\begin{center}
\begin{tabular}{lcccc}
\toprule
任务类型 & 腿数 & 行程 m & 占比 & 折算 s/源 \\
\midrule
普查站位 & 22.8 & 13246 & 60.7\% & 212 \\
清除 & 26.0 & 6233 & 28.6\% & 100 \\
补测站位 & 5.9 & 2336 & 10.7\% & 37 \\
\midrule
行程合计 & & 21815 & & 349 \\
检测 424 次 & & & & 202 \\
\midrule
总计 & & & & 551 \\
\bottomrule
\end{tabular}
\end{center}

实测 595.9 秒每源与该分解基本一致，说明时间构成中没有明显浪费。

三项关键测量确定了压缩空间：

第一，实测路径 21583 m，而对同一批点的 2-opt 巡回是 19964 m，
余量只有 8\%，路由已近最优。

第二，清除任务行程 8633 m，而 13 个源坐标的 2-opt 巡回是 8635 m，
两者几乎完全相同，已到下界。

第三，检测次数在 8 组参数下稳定在 420 至 440 次之间，已到自身下界。

% =====================================================================
\section{不可达目标的界定}

原始目标是把问题四压缩到 300 至 400 秒每源。经核算，该目标在当前问题设定下不可达，
依据如下。

一局必须访问约 58 个位置点，其中普查站位约 19 个、源 12.5 个、
归航中间点约 26 段。把这批点用 2-opt 求最优巡回约 20000 m，
按 5 m/s 折合 4000 秒，即 320 秒每源，尚未计入测量就已经逼近目标区间下界。

减少点数的三条路径都被实测堵死。减少普查站位会触发自适应补站，
而自适应站位是逐步贪心选择的，没有全局视野，路线质量明显更差：
站位数由 24.1 降到 14.6，行程反而由 24010 m 升到 27513 m。
减少贴边环会使贴边源的发现率下降。减少归航中间点会牺牲完成率。

因此本文把 300 至 400 秒记为在当前设定下不可达，
现实下界在 450 秒每源以上。若要真正逼近该区间，需要改变问题本身的一项，
例如提高有效接收半径、允许一次测量覆盖多个信道，或把严格证明改为置信度判据。

% =====================================================================
\section{负结果与经验}

研究过程中实测了 9 项结构性替代方案与 30 余项参数尝试，其中大部分为负结果。
完整清单见随附的尝试日志，此处归纳三条反复出现的规律。

第一，靠增加预算或扩大图案来补救，几乎总是负优化。
提高清除尝试次数、加大末段半径、扩大补测范围，都会增加时间而不改善完成率。
需要的是换个位置观测，而不是在原地更努力。

第二，采样密度存在最优点。格点间距过密或过疏都变差，950 m 附近是甜点。
认证判据给出了这个甜点的几何解释。

第三，显式实现的耦合不一定带来增量。任务排序中已有的优先权偏置，
往往已经隐含实现了设计者想要引入的耦合。

此外，项目过程中一共出现五次单一算例池误导结论的情况，
因此最终选型一律采用四池联合、以最差池为准。

% =====================================================================
\section{局限与后续方向}

完成率的剩余缺口集中在贴边朝外的定向源。这类源在域内的受光区是极窄的月牙，
需要沿边界密集布站才能保证覆盖，时间预算不允许。
若要改进，方向是主动重捕获：一旦某频道被听到但无法定位，
就回到该听到点，沿垂直于该方位的方向做一系列短步探测，
步长由信号是否仍然存在二分确定，直到取得方向差足够的第二条方位。

时间方面，剩余的可压缩项是清除任务的接近行程与归测过程。
若要显著下降，需要把认证站位与源的位置联合规划，
使两类目标共享行程，而不是各自独立布点。

\end{document}
